\documentclass[12pt,reqno]{amsart}
\usepackage{amsmath,amssymb,amsthm}
\usepackage{url}
\usepackage{xspace}
\usepackage[hidelinks]{hyperref}
\usepackage{xcolor}
\usepackage{enumerate}
\usepackage{verbatim}
\usepackage{array}
%\usepackage{blindtext}
%\usepackage{scrextend}
%\addtokomafont{labelinglabel}{sffamily}
%\usepackage[symbol]{footmisc}
%\renewcommand*\rmdefault{cmfib}
\renewcommand{\thefootnote}{\ensuremath{\fnsymbol{footnote}}}


%\usepackage[light,math]{iwona}
%\usepackage[finnish]{babel}
\newcounter{prob}
\setcounter{prob}{1}
\newcounter{assign}
\setcounter{assign}{1}
\newcounter{spacerule}
%\renewcommand{\theprob}{\Roman{prob}}
%\newfont{\bigrm}{rm at 12pt}

\newenvironment{problem}{\par\noindent\textbf{Problem \Roman{prob}.\hskip2pt}}{\par\vskip12pt\stepcounter{prob}}
\newenvironment{assignment}{\par\noindent\textbf{Assignment \Roman{assign}.\hskip2pt}}{\par\vskip15pt\stepcounter{assign}}
\usepackage{color}
\definecolor{heading}{RGB}{40, 0, 130}
\definecolor{bgcolor}{RGB}{255,253,250}
\definecolor{duecolor}{RGB}{150,0,0}
\definecolor{hyperlink}{RGB}{33,0,255}
\pagecolor{bgcolor}


%-----------------------------------------------------------------------------%
% PAGE SIZES
%-----------------------------------------------------------------------------%
\setlength{\hfuzz}{3pt}
\setlength{\headheight}{32pt}
\setlength{\headsep}{29pt}
\setlength{\footskip}{28pt}
\setlength{\textwidth}{444pt}
\setlength{\textheight}{636pt}
\setlength{\marginparsep}{7pt}
\setlength{\marginparpush}{7pt}
\setlength{\oddsidemargin}{4.5pt}
\setlength{\marginparwidth}{55pt}
\setlength{\evensidemargin}{4.5pt}
\setlength{\topmargin}{-15pt}
\setlength{\footnotesep}{8.4pt}
%--------------------------------------------------------------------%
% CLAIMS
%--------------------------------------------------------------------%
%\swapnumbers
\theoremstyle{plain}
\newtheorem{theorem}{Theorem}
\newtheorem{corollary}[theorem]{Corollary}
\newtheorem{lemma}[theorem]{Lemma}
\newtheorem{proposition}[theorem]{Proposition}
%
\theoremstyle{definition}
\newtheorem{remark}[theorem]{Remark}
\newtheorem{definition}[theorem]{Definition}
\newtheorem{example}[theorem]{Example}
%--------------------------------------------------------------------%
% OPERATORS
%--------------------------------------------------------------------%
\DeclareMathOperator{\pr}{pr}
\DeclareMathOperator{\tot}{tot}
\let\hook\undefined
\DeclareMathOperator{\hook}{\lrcorner}
\DeclareMathOperator{\Sym}{Sym}
\DeclareMathOperator{\Div}{Div}
\DeclareMathOperator{\D}{D}
\DeclareMathOperator{\tr}{tr}
%
%     Macros
%
\newcommand\ianhook{\mathbin{\raise1pt\hbox{\hbox{{\vbox{\hrule height.4pt
width6pt depth0pt}}}\vrule height3pt width.4pt depth0pt}\,}}
\newcommand{\volform}{dx^1 dx^2\cdots dx^\bdim}
\newcommand{\R}[1]{\mathbb{R}^{#1}}
\newcommand{\M}[1]{\mathbb{M}^{#1}}
\newcommand{\krondel}[2]{\delta^{#1}_{#2}}
\newcommand{\He}{Helmholtz\xspace}
\newcommand{\ga}[1]{\mathfrak{ga}(#1)}
\newcommand{\cprime}{\/{\mathsurround=0pt$'$}}
\newcommand{\comp}{\raise 1pt \hbox{$\,\scriptstyle\circ\,$}}


\newcommand{\cB}{\mathcal{B}}
\newcommand{\cL}{\mathcal{L}}
\newcommand{\cE}{\mathcal{E}}
%\newcommand{\BN}{\mathbb{N}}
\newcommand\closure[2]{\overset{\raisebox{0.6pt}{\hspace{0.2em}\text{\rule{#1}{0.2mm}}}}{#2}}
\newcommand{\curve}{\boldsymbol{\alpha}}
\newcommand{\curveb}{\boldsymbol{\beta}}
\newcommand{\curvature}{\boldsymbol{\kappa}}
\newcommand{\torsion}{\boldsymbol{\tau}}
\newcommand{\T}{\mathbf{T}}
\newcommand{\PN}{\mathbf{N}}
\newcommand{\BN}{\mathbf{B}}
\newcommand{\unitv}{\mathbf{u}}
\newcommand\plane{\mathcal{P}}
%\newcommand\ianhook{\mathbin{\raise1pt\hbox{\hbox{{\vbox{\hrule height.4pt width6pt depth0pt}}}\vrule height3pt width.4pt depth0pt}\,}}
%\newcommand{\ianhook}{\lrcorner}
\newcommand\bdim{m}
\newcommand\bcoord[1]{x^{#1}}
\newcommand\Christoffel[3]{\Gamma^{#1}_{#2#3}}
\newcommand\Christoffelvar[4]{\Gamma^{\hphantom{#1}#2}_{\mkern-3mu#1,#3#4}}
\newcommand\CP[1]{\mathbb C\text{P}^{#1}}
%\newcommand\curvature[4]{R_{#1#2#3}{}^{#4}}
\newcommand\cx[1]{\mathbb C^{#1}}
\newcommand\cxp[1]{\mathbb C^{#1^\mathbf{*}}}
\newcommand\dvChristoffel[3]{\gamma^{#1}_{#2#3}}
\newcommand\func{F}
\newcommand\funcg{G}
\newcommand\gc[2]{g_{#1#2}}%components of the metric
\newcommand\gcvar[3]{g_{#1,#2#3}}
\newcommand\gci[2]{g^{#1#2}}%components of the inverse
\newcommand{\gcivar}[3]{g_{#1}^{#2#3}}
\newcommand\gcj[3]{g_{#1#2,#3}}%jets of the metric
\newcommand\gcjvar[4]{g_{#1,#2#3,#4}}
\newcommand\gccoll[1]{g^{[#1]}}
\newcommand\gdet{g}
\newcommand{\gdetvar}[1]{g_{#1}}
\newcommand\gvol{\nu_g}
\newcommand\gvolform{\gdet^{1/2}d\bcoord 1\wedge\cdots\wedge d\bcoord\bdim}
\newcommand\gind{l}
\newcommand{\gindd}{p}
\newcommand\Gbundle{\mbox{\sffamily{G}}}
\newcommand\Gr[2]{\mathbf{Gr}(#1,#2)}
\newcommand\Grs[2]{\mathbf{G}_{#1#2}}
\newcommand\RP[1]{\text{P}^{#1}}
\newcommand\pnt{p}
\newcommand\projP{\pi_{\text{P}}}
\newcommand\projGr{\pi_{{Gr}}}
\newcommand\rktwos{\mathcal T}
\newcommand\rktwo[2]{\rktwos(#1,#2)}
\newcommand\smfld{S}
\newcommand\ib{\mathbf i}
\newcommand\jb{\mathbf j}
\newcommand\kb{\mathbf k}
\newcommand\rb{\mathbf r}
\newcommand\ub{\mathbf u}
\newcommand\vb{\mathbf v}
\newcommand\wb{\mathbf w}
\newcommand\vsp{\mathcal V}

\newcommand\on[1]{\mathbf e_{#1}}
\newcommand\GrEcoord[2]{\varphi^e_{#1#2}}
\newcommand\Grcoord[2]{\varphi_{#1#2}}
\newcommand\Grcoordinv[2]{\psi_{#1#2}}
\newcommand\Emx{\mathbb E}
\newcommand\Transpose[1]{{#1}^{\text{\bf T}}}
\newcommand\ont[1]{\mathbf e^{\text{\bf T}}_{#1}}
\newcommand\mx[2]{\mathcal M_{#1\times #2}}
\newcommand\Pcoord[1]{\varphi_{#1}}
\newcommand\openset{\mathcal U}
\newcommand\vf{X}
\newcommand\vfalt{Y}

\def\QED{\hskip0.1em\hfill\null\
  \null\nobreak\hfill\kern3pt\vbox{\hrule\hbox
   {\vrule\kern1pt\vbox{\kern1.7pt\hbox{$\scriptscriptstyle{QED}$}
    \kern0.2pt}\kern1pt\vrule}\hrule}}
\renewcommand{\qedsymbol}{\QED}

\def\END{\hskip0.1em\hfill\null\
  \null\nobreak\hfill\kern3pt\vbox{\hrule\hbox
   {\vrule\kern1pt\vbox{\kern1.7pt\hbox{$\,\,\,\vspace{5pt}$}
    \kern0.2pt}\kern1pt\vrule}\hrule}}

\newcommand{\abs}[1]{\lvert\kern.1em{#1}\rvert}
\newcommand{\norm}[1]{\lVert\kern.1em{#1}\rVert}
\newcommand*{\pd}[2]{\mathchoice{\frac{\partial #1}{\partial #2}}
  {\partial #1/\partial #2}{\partial #1/\partial #2} {\partial
  #1/\partial #2}}
\newcommand*{\od}[2]{\mathchoice{\frac{d#1}{d#2}}
  {d#1/d#2}{d#1/d#2}{d#1/d#2}}

\hyphenation{Noeth-er Helm-holtz Poh-jan-pel-to}
\newcommand\lasku{\count1=1}
\newcommand\katso{\showthe\count1}
\begin{document}
\centerline{\large{\scshape{\color{heading} MTH 430 Metric Spaces and Topology}}}
\vskip3pt
\centerline{\large{\scshape{\color{heading} Spring Term 2021}}}
\vskip3pt
\centerline{\large{\ttfamily Midterm Exam}}
\vskip2pt
%\centerline{\em{\color{duecolor} Due April 9, 2019}}
\vskip15pt

%\begin{problem}
%Show that in a Hausdorff space, any point is a closed set.
%\end{problem}
%\begin{problem}
%Show that the real line is homeomorphic with the interval $(0,1)$.
%\end{problem}
%\begin{problem}
%Explain why any closed subset of a compact space must be compact.
%\end{problem}
\begin{problem}
Let $f\colon X\to Y$ be a function where $X$ is a topological space. Let $\tau_{Y,f}=\{U\subset Y\,\vert\,\mbox{$f^{-1}(U)\subset X$ is open}\}$. Show that $\tau_{Y,f}$ is a topology on $Y$.
\end{problem}

\begin{problem} Let $X$ be a topological space and let $A\subset X$. Show that the closure $\closure{8pt}{A}$ (the union of A with all its limit points) of $A$ is the smallest closed set containing $A$. 
\end{problem}

\begin{problem} Let $X$ be a Hausdorff space, and let $\mathcal C\subset X$ be compact. Show that for every $p\in X\setminus \mathcal C$ there are disjoint open sets $\mathcal O_1$, $\mathcal O_2$ such that $\mathcal C\subset \mathcal O_1$ and $p\in\mathcal O_2$. 
\end{problem}
\begin{problem}
\setcounter{spacerule}{1}
\begin{list}{\textbf{\alph{spacerule})}}{\setlength{\itemsep}{2pt}\setlength{\leftmargin}{30pt}\setlength{\parsep}{0pt}\usecounter{spacerule}}
\item Show that $f\colon X\to Y$ is continuous if and only if $f^{-1}(C)\subset X$ is closed for every closed $C\subset Y$.
\item Then show that a function $f\colon X\to Y$ is continuous if and only if $f(\overset{\raisebox{1pt}{\hspace{0.2em}\text{\rule{7pt}{0.2mm}}}}{A})\subset \overset{\raisebox{1pt}{\hspace{0.1em}\text{\rule{20pt}{0.2mm}}}}{f(A)}$ for all $A\subset X$.
\end{list}
\end{problem}
%\begin{problem}
%Let $f\:S\to T$ be continuous (in the $\epsilon$-$\delta$-sense), where $S$, $T$ are 
%metric spaces. Show that the inverse image $f^{-1}(\mathcal O)$ of an open set $\mathcal O\subset T$
%is open in $S$. 
%\end{problem}
%\begin{problem} Suppose $Z\subset X$ is a connected subset of a topological space $X$ in the relative topology. Show that the closure $\overline Z$ is also connected.
%\item{3.} Show that the closure $\overline{A} = A\cup A'$, where $A'$ is the set of
%accumulation points of a set $A\subset X$ in a topological space $X$.
%\end{problem}
\begin{problem}
Let $\beta=\{[a,b)\;|\; -\infty<a<b<\infty\}$ be the collection
of half-closed finite intervals on $\Bbb R$, and let $\tau$ be the topology generated by it. Show that $\beta$ is indeed a basis for $\tau$.
Is $\tau$ Hausdorff?
Is $\tau$ coarser or finer (or neither) than the 
standard topology of $\R{}$?
\end{problem}

%\begin{problem}
%Let $f$ be a continuous bijection from a compact space $X$ onto a Hausdorff space $Y$. Show that $f$ is a homeomorphism.
%\end{problem}

\begin{problem} Let $X$ be a topological space and let $A\subset X$ be both open and closed. Suppose $B\subset X$ is connected. Show that either $B\subset A$ or $A\cap B=\emptyset$.

\end{problem}
 
\begin{problem} Let $K_1\supset K_2\supset K_3\supset\cdots$ be a sequence of non-empty, closed and compact sets in some topological space $X$. Show that $K_1\cap K_2\cap K_3\cap\cdots\neq\emptyset$.
\end{problem}
%\begin{problem}
%Let $(X,d)$ be a metric space. Show that $x\in X$ is in the closure of a set $A\subset X$ if and only if $\inf\{d(x,a)\,\vert\,a\in A\}=0$.
%\end{problem}
%\begin{problem}
%Show that the the finite complements topology on any set $X$ is the coarsest topology in which single points are closed sets.
%\end{problem}
\end{document}
\vskip5pt 
\begin{problem} Recall that a function is continuous if the inverse image of any open set in $Y$ is open in $X$.
\begin{enumerate}[(a)]
\item Since any open set in the original topology is also open in any finer topology of $X$, $f$ is continuous in the new topology.
\item Let $\tau_{X,f}=\{f^{-1}(U)\subset X\,\vert\,\mbox{$U\subset Y$ is open}\}$ be a collections of subsets of $X$. It is easy to see that $\tau_{X,f}$ is a topology, cf.~assignment I, problem 2, and by construction, it is the coarsest topology on $X$ in which $f$ is continuous. Thus, given a topology $\tau_1$ on $X$, $f$ is continuous if $\tau_1$ is finer than $\tau_{X,f}$ and $f$ is not continuous if $\tau_1$ is properly coarser than $\tau_{X,f}$.
\item Let $\tau_{Y,f}=\{U\subset Y\,\vert\,\mbox{$f^{-1}(U)\subset X$ is open}\}$. Again one can check that $\tau_{Y,f}$ is a topology, and by definition, it is the finest topology on $Y$ in which $f$ is continuous. Thus given a topology $\tau_2$ on $Y$, $f$ is continuous if $\tau_2$ is coarser than $\tau_{Y,f}$ (cf.~part (d)) and $f$ is not continuous if $\tau_2$ is properly finer than $\tau_{Y,f}$.  
\item Since any set open in the new topology is also open in the original topology, $f$ is continuous in the new topology.
\end{enumerate}
\item
Let $X$ and $Y$ be two topological spaces, and let $f\colon X\to Y$ be some map.
Equip the image $f(X)\subset Y$ with the relative topology. Show that $f$ is continuous if and only if $f\colon X\to f(X)$ (that is, $f$ considered as a map onto its image) is continuous. 
\vskip5pt 
\noindent
{\bf Solution:} Write $\widehat{f}$ for the map onto $f(X)$. Suppose first that $f\colon X\to Y$ is continuous, and let $U\subset f(X)$ be open. Then $U=f(X)\cap O$, where $O\subset Y$ is open. Now $\widehat{f}^{-1}(U)=f^{-1}(O)$ (check!), which is open by the continuity of $f$.
\vskip3pt
\noindent
Next assume that $\widehat{f}$ is continuous, and let $O\subset Y$ be open. Then $f^{-1}(O)=\widehat{f}^{-1}(O\cap f(X))$, which is open as $O\cap f(X)\subset f(X)$ is open in the relative topology.  
\item 
\begin{enumerate}[(a)]
\item Let $X$ a topological space, and suppose that there is $x\in X$ such that $X\setminus\{x\}$ is a union of two disjoint open sets. Let $Y$ be homeomorphic with $X$. Show that there is $y\in Y$ such that $Y\setminus\{y\}$ is a union of two disjoint open sets.
\item Show that the real line $\R{}$ and the plane $\R2$ can not be homeomorphic.
\end{enumerate}

\newpage
\noindent
{\bf Solution:}
\begin{enumerate}[(a)]
\item Let $f\colon X\to Y$ be the homeomorphism, and let $y=f(x)$. Then $f\colon X\setminus\{x\}\to Y\setminus\{y\}$ is a bijection. Since $f$ is a homeomorphism, it is also an open map. Write $X\setminus\{x\}=U_1\cup U_2$, where $U_1$, $U_2$ are open and disjoint. Then $Y\setminus\{y\}=f(U_1)\cup f(U_2)$, where $f(U_1)$, $f(U_2)$ are open and disjoint.
\item Suppose $\R{}$ and $\R2$ were homeomorphic under $f$. Now $\R{}\setminus\{0\}$ is the union of two disjoint open sets, so by part (a), $\R2\setminus\{f(0)\}$ would also be a union of two disjoint open sets. But this is a contradiction as any two points in $\R2\setminus\{f(0)\}$ can be connected with a continuous curve. (We will return to this when we discuss path connectedness.) 
\end{enumerate}
\item 
\begin{enumerate}[(a)]
\item Let $X$ be any set, and let $\tau=\{A\subset X\,\vert\,\mbox{ $A=\emptyset$ or $X\setminus A$ is finite}\}$. Show that the collection $\tau$ of subsets of $X$ forms a topology. (This is so-called finite complements topology.)
\item Show that the finite complements topology on the real axis is not Hausdorff.
\end{enumerate}
\vskip5pt 
\noindent
{\bf Solution:}
\begin{enumerate}
\item Clearly $\emptyset$, $X\in\tau$. Let $U_\alpha$, $\alpha\in\mathcal A$, be any sets in $\tau$. Then $(\cup_{\alpha\in\mathcal A}U_\alpha)^\complement=\cap_{\alpha\in\mathcal A}U_\alpha^\complement$, which is finite, so $\cup_{\alpha\in\mathcal A}U_\alpha\in\tau$. Here $V^\complement$ denotes the complement of a set $V$. Next let $U_1$,\dots, $U_n\in\tau$. Then $(U_1\cap\cdots\cap U_n)^\complement = U_1^\complement\cup\cdots\cup U_n^\complement$, which again is finite. Thus $U_1\cap\cdots\cap U_n\in\tau$.
\item Consider $0$, $1\in\R{}$. Let $U_o$, $U_1$ be any two neighborhoods of $0$, and $1$, respectively. Then $(U_o\cap U_1)^\complement = U_o^\complement\cup U_1^\complement$, which is finite. As there are infinitely many real numbers, $U_o\cap U_1\neq\emptyset$. 
\end{enumerate}
\item 
\begin{enumerate}[(a)]
\item Let $X$ be a compact space, and let $A\subset X$ be closed. Show that $A$ is compact in the relative topology.
\item Show by the way of an example that the above claim may fail if $A$ is not closed in $X$.
\end{enumerate}
\vskip5pt 
\noindent
{\bf Solution:}
\begin{enumerate}
\item Let $U_\alpha$, $\alpha\in\mathcal A$, be an open cover of $A$. By definition, $U_\alpha=A\cap O_\alpha$, where $O_\alpha\subset X$ is open. Now the sets $O_\alpha$, $\alpha\in\mathcal A$, $A^\complement$, form an open cover of $X$. But $X$ is compact, so it admits a finite subcover $O_{\alpha_1}$,\dots,$O_{\alpha_n}$, $A^\complement$. Then $U_{\alpha_1}$,\dots,$U_{\alpha_n}$ covers $A$ so $A$ must be compact.
\item By Heine-Borel, the closed interval $[0,1]$ is compact, but the open interval $(0,1)\subset[0,1]$, as we've seen in class, is not. 
%$\raisebox{-5pt}{\Large\urcorner}$
\end{enumerate}
\end{list}
\vskip10pt
\noindent
{\bf Extra: Solution to problem 1(c), assignment II:} Recall that a set $O\subset\R{}$ is open if for any $x\in O$, there is an $\epsilon>0$ such that the open interval $(x-\epsilon,x+\epsilon)\subset O$. Thus, given the definition of an accumulation point, the rationals $\mathbb Q{}$ are dense in $\R{}$, if, given any real number $r\in\R{}\setminus\mathbb Q$ and $\epsilon>0$, there is some $q\in\mathbb Q$ such that $|r-q|<\epsilon$. This can be seen hold true by truncating the decimal expansions of real numbers. As an example, which will also show how to do the general case, consider $\pi=3.14159265359\dots$, and let $\epsilon_1=0.1$. Choose $q_1=3.1=31/10$. Then $|\pi-q_1|=0.041\dots<\epsilon_1$. Next, say, let $\epsilon_2=10^{-5}$. Choose $q_2=3.14159=314159/100000$. Then $|\pi-q_2|=0.00000265\dots<10^{-5}$. And so on. 
\end{document}
\vskip15pt
\noindent
\begin{problem}
\begin{list}{\textbf{\alph{spacerule})}}{\setlength{\itemsep}{4pt}\setlength{\leftmargin}{30pt}\setlength{\parsep}{0pt}\usecounter{spacerule}}
\item If sets $C$ and $D$ satisfy $C\subset D$ and $D\subset C$, then $C=D$. Thus, pick $y\in f(A_1\cup A_2)$ so that $y=f(a)$ for some $a$ contained in $A_1$ or $A_2$ (or both). Now it follows that $y\in f(A_1)\cup f(A_2)$. And so on \dots. 
%\item Again, choose an arbitrary element in the set on the left-hand side and show that it must be contained in the set on the right-hand side. It is also easy to find examples in which the inclusion is strict. 
\end{list}
\end{problem}
\begin{problem}
One can derive these identities as in problem \#1.
\end{problem}
\begin{problem}
The problem can be dealt with by proving that $(a)\Rightarrow(b)\Rightarrow(c)\Rightarrow(d)\Rightarrow(e)\Rightarrow(a)$.

\vspace{3pt}
\begin{tabular}{lp{13.4cm}}
$(a)\Rightarrow(b):$ & Suppose not. Since $A\subset f^{-1}(f(A))$ for any $A$, there is some\vfil ${b\in f^{-1}(f(A))\setminus A}$. Then $f(b)\in f(A)$, so $f(b)=f(a)$ for some $a\in A$. Thus $f$ can not be injective.\\
$(b)\Rightarrow(c):$ & Again, suppose not. Already know that $f(A\cap B)\subset f(A)\cap f(B)$, so there must be some $y\in f(A)\cap f(B)\setminus f(A\cap B)$. Then $y=f(a)$, $y=f(b)$ for some $a\in A\setminus A\cap B$, $b\in B\setminus A\cap B$. Thus $\{a\}\subsetneq f^{-1}(f(\{a\})$ so (b) fails.
\\
$(c)\Rightarrow(d):$ & Obvious.\\
$(d)\Rightarrow(e):$ & Now $A=(A\setminus B)\cup B$, so $f(A) = f(A\setminus B)\cup f(B)$, and, by (d), $f(A\setminus B)\cap f(B)=\emptyset$. Thus $f(A\setminus B)=f(A)\setminus f(B)$.\\
$(d)\Rightarrow(e):$ & Suppose that $f(a)=f(b)=y$ for some $a\neq b$. Let $A=\{a,b\}$, $B=\{b\}$. Then $f(A\setminus B)= \{y\}$ but $f(A)\setminus f(B)=\emptyset$. Thus (e) fails.
\end{tabular} 
\end{problem}
\begin{problem}
Obviously, $\emptyset, X\in\tau$, and $\tau$ is closed under unions. Hence we only need to prove that $\tau$ is closed under finite intersections. First, it follows from the second assumption that the intersection $B_1\cap B_2\in\tau$ for any $B_1, B_2\in\beta$.

Next, given any $X, Y\in\tau$, write 
\[X=\cup_{\kappa\in\mathcal K} B_{1,\kappa},\quad Y=\cup_{\lambda\in\mathcal L} B_{2,\lambda}, \quad\text{where $B_{1,\kappa}, B_{2,\lambda}\in\beta$}.
\] 
Then, by the above, 
\[
X\cap Y = \cup_{\kappa\in\mathcal K, \lambda\in\mathcal L} B_{1,\kappa}\cap B_{2,\lambda}\in\tau.
\] 
Finally, use induction in the number of sets in an intersection. 
\end{problem}
\begin{problem}
The problem simply boils down to finding inverse images of elements of $\beta$ under the given mapping. Note that the topology $\tau$ determined by $\beta$ consists of unions of intervals $[a,b)$, $a<b$. Thus, for example, in part (b), $f^{-1}([0,1))=(-1,0]$, which is not contained in $\tau$. Thus $f$ can not be continuous.   
\end{problem}
\end{document}
In this problem you need to show that the collection of sets $\tau=\{\text{all unions of $B\in\beta$}\}$ is a topology, that is, it satisfies,
\begin{list}{\textbf{\arabic{spacerule})}}{\setlength{\itemsep}{4pt}\setlength{\leftmargin}{30pt}\setlength{\parsep}{0pt}\usecounter{spacerule}}
\item $\emptyset$, $X\in\beta$. 
\item If $\{B_\alpha\}$ is any collection of sets in $\beta$, then the union $\cup_\alpha B_\alpha\in\tau$.
\item If $B_1$, $B_2$, \dots, $B_k\in\tau$, then the finite union $B_1\cap B_2\cap\cdots\cap B_k\in\tau$ for any $B_1, B_2\in\beta$. Next, given any $X, Y\in\tau$, write 
\[X=\cup_{\kappa\in\mathcal K} B_{1,\kappa},\quad Y=\cup_{\lambda\in\mathcal L} B_{1,\lambda}
\] 
\end{list} 
Conditions 1 and 2 are more or less immediate. As to 3, first show that the intersection $B_1\cap B_2$ of two sets $B_1$, $B_2\in\tau$ again lies in $\tau$. 
\newcounter{nested}
\begin{assignment}({\em{\color{duecolor} Due Thursday, April 9th, 2020}})
\begin{list}{\textbf{\arabic{spacerule})}}{\setlength{\itemsep}{4pt}\setlength{\leftmargin}{30pt}\setlength{\parsep}{0pt}\usecounter{spacerule}}
\item Let $f\colon X\to Y$ be a function. If $A\subset X$ and $B\subset Y$, then the image $f(A)$ of $A$ and the inverse image $f^{-1}(B)$ of $B$ are defined by 
$$f(A)= \{f(a)\,\vert\, a\in A\},\quad\quad f^{-1}(B) = \{a\in A\,\vert\, f(a)\in B\}.$$
Prove the following statements.
\begin{enumerate}[(a)]
\item For all $A_1$, $A_2\subset X$, $f(A_1\cup A_2)=f(A_1)\cup f(A_2)$.
\item For all $A_1$, $A_2\subset X$, $f(A_1\cap A_2)\subset f(A_1)\cap f(A_2)$.
\end{enumerate}

\item Prove the following statements.
\begin{enumerate}[(a)]
\item For all $B_1$, $B_2\subset X$, $f^{-1}(B_1\cup B_2)=f^{-1}(B_1)\cup f^{-1}(B_2)$.
\item For all $B_1$, $B_2\subset X$, $f^{-1}(B_1\cap B_2)=f^{-1}(B_1)\cap f^{-1}(B_2)$.
\end{enumerate}
\item (15 points) Let $f\colon X\to Y$ be a function. Prove that the following are equivalent.
\begin{enumerate}[(a)]
\item $f$ is injective.
\item For all $A\subset X$, $f^{-1}(f(A))=A$.
\item For all $A$, $B\subset X$, $f(A\cap B)=f(A)\cap f(B)$.
\item For all $A$, $B\subset X$, $A\cap B=\emptyset$ implies that $f(A)\cap f(B)=\emptyset$.
\item For all $A$, $B\subset X$ with $B\subset A$, $f(A\setminus B)=f(A)\setminus f(B)$.
\end{enumerate}
\item Let $X$ be a set and let $\beta$ be a collection of subsets of $X$ such that 
\begin{enumerate}[(a)]
\item For every $p\in X$, there is some $B\in\beta$ such that $p\in B$.
\item Let $p\in B_1\cap B_2$, where $B_1, B_2\in\beta$. Then there is some $B_3\in\beta$ such that $p\in B_3\subset B_1\cap B_2$.
\end{enumerate}
Show that the family of subsets of $X$ consisting of all unions of the members of $\beta$ together with the empty set $\emptyset$ form a topology on $X$. Explain why the above statement is stronger (and hence more useful) that Theorem 2.5 in the course text.
\item Let $X=\R{}$ and suppose that $\R{}$ is equipped with the topology defined by the basis $\beta=\{[a,b)\,\vert\,a<b\}$. Call this space $\R{}_\tau$. Which of the following functions $f\colon \R{}_\tau\to\R{}_\tau$ are continuous?
\begin{enumerate}[(a)]
\item $f(x) = 2x-5$
\item $f(x) = -x$
\item $f(x) = x^2$
\end{enumerate}
\end{list}
\end{assignment}
\end{document}

\item Let $\mathcal P=\{(x,y,x,w)\in\R4\,\vert\, x,y,w\in \R{}\}$ and $\mathcal Q=\{(x,y,z,x)\in\R4\,\vert\,x,y,z\in\R{}\}$. Describe the intersection $\mathcal P\cap\mathcal Q$.
\item Find all vectors $\vb$ in $\R5$ that is orthogonal to all the three vectors
$$
<1,2,0,0,-2>,\quad <-2,1,2,1,0>,\quad <0,-2,1,0,2>.
$$
\item Let $R_\alpha$ denote the rotation in $\R3$ about the axis determined by the vector $\vb = \jb + 2\kb$ by angle $\alpha$. (Here the $\kb$ component of the image of $\jb$ is positive for small $\alpha$.) Find the images of $\ib$, $\jb$, and $\kb$ under $R_\alpha$.
\item Show that $\ub\times(\vb\times\wb) = (\ub\times\vb)\times\wb +\vb\times(\ub\times\wb)$ for any vectors $\ub$, $\vb$, $\wb$. Thus the cross product is not associative.
\item A fluid is flowing at a constant velocity of $<1,-1,0>$. Find the volume of fluid flowing through the triangle determined by $<1,0,1>$ and $<2,-1,1>$ in one second. (Here and elsewhere in this course we use metric units.)  
\item Strang, \S 11.3, problems 37, 47, 48.
\item OpenStax, Ch.~2, problems 70, 107, 219, 283. 
\end{list}
\end{assignment}
\begin{assignment}({\em{\color{duecolor} For the week ending on January 22, 2020}})
\begin{list}{\textbf{\arabic{spacerule})}}{\setlength{\itemsep}{4pt}\setlength{\leftmargin}{30pt}\setlength{\parsep}{0pt}\usecounter{spacerule}}
\item Let $f(x,y)=x^2y$, and let $D=[0,1]\times[0,1]$ be the unit square. Suppose $D$ is partitioned by $\mathcal P_n$, $n\in\mathbb N$, where $\mathcal P_n$ consists of the squares $[\dfrac{k-1}{2^n},\dfrac{k}{2^n}]\times[\dfrac{l-1}{2^n},\dfrac{l}{2^n}]$, $1\leq k,l\leq 2^n$. Find the upper and lower sums $U(f,\mathcal P_n)$, $L(f,\mathcal P_n)$ for $n\in\mathbb N$, and show that $f(x,y)$ is integrable over $D$. What is the value of the integral $\iint_D f\,dA$?
\item Let $F(x,y) = (2x-y,5x+3y)$. Find the inverse function for $F(x,y)$ and then compute the Jacobians of $F(x,y)$ and its inverse. What do you notice?
\item 
Find the Jacobian of
\begin{enumerate}[(a)]  
\item  $F(x,y) = (\text{e}^x\cos y,\text{e}^x\sin y)$. 
\item $F(x,y)=(f(y), g(x))$. When is $F(x,y)$ orientation preserving/reversing?
\end{enumerate}
\item Solve problems 120 and 121 on page 525 in OpenStax using elementary Euclidean geometry.
\item OpenStax Ch.~5, problems 86, 89, 97, 98, 111, 113, 361, 363, 370, 388. 
\item Strang, \S 14.2, problems 17, 19, 21.
\end{list}
\end{assignment}
\begin{assignment}({\em{\color{duecolor} For the week ending on January 29, 2020}})
\begin{list}{\textbf{\arabic{spacerule})}}{\setlength{\itemsep}{4pt}\setlength{\leftmargin}{30pt}\setlength{\parsep}{0pt}\usecounter{spacerule}}
\item Let $T$ be the transformation from the $uv$-plane into the $xy$-plane defined by $x=u^2-v^2$, $y=2uv$. 
\begin{enumerate}[(a)]
\item Find the image under $T$ of the unit square $S=\{(u,v)\,\vert\,0\leq u\leq1, 0\leq v\leq1\}$ in the $xy$-plane.
\item Introduce polar coordinates $r$, $\theta$ in the $uv$-plane. Find the image under $T$ of the wedge $0<r<2$, $0<\theta<\dfrac\pi{3}$ in the $xy$-plane.
\end{enumerate}

\item Let $S$ be the region in the $xy$-plane bounded by the lines $y-x=0$, $y-x=2$, $y+x=0$, $y+x=2$. 
\begin{enumerate}[(a)]
\item Find a transformation from the $xy$-plane into the $uv$-plane such that the image of $S$ is the $uv$-plane is the unit square $\{(u,v)\,\vert\,0\leq u\leq1, 0\leq v\leq1\}$. 
\item Evaluate the integral $\iint_S \sqrt{y^2-x^2}\, dA$ by a change of variables under the transformation you found in part a. 
\end{enumerate}
\item Let $S$ be the the trapezoid bounded by the lines $y=x+1$, $y=x+3$, $y=4-2x$, $y=5-x$. Find a transformation from the $xy$-plane into the $uv$-plane such that the image of $S$ is the unit square $\{(u,v)\,\vert\,0\leq u\leq1, 0\leq v\leq1\}$. 
\item Let $D_\lambda\colon\R2\to\R2$, $\lambda>0$, be the dilation $D_\lambda(x,y)=(\lambda x,\lambda y)$. Suppose the a region $S$ has area $A(S)$. Find the area of $D_\lambda(S)=\{(\lambda x,\lambda y)\,\vert\,(x,y)\in S\}$.
\item  Denote $\R2$ with the origin removed by $\R{2^*}$. Let $\text{Inv}\colon\R{2^*}\to\R{2^*}$ be the inversion $\text{Inv}(x,y)=(\dfrac{x}{x^2+y^2},\dfrac{y}{x^2+y^2})$. Find the inverse function and the Jacobian of $\text{Inv}$. What is the image of the circle $x^2+y^2=r^2$ under $\text{Inv}$?
\item OpenStax Ch.~5, problems 389, 397, 400 (fix a typo!), 402.
\item OpenStax Ch.~6, problems 56, 59, 64 (solve these by hand). Fix a typo in example 6.19.
\item Strang, \S 15.2, problems 6, 19, 21. 
\end{list}
\end{assignment}
\begin{assignment}({\em{\color{duecolor} For the week ending on February 12, 2020}})
\begin{list}{\textbf{\arabic{spacerule})}}{\setlength{\itemsep}{4pt}\setlength{\leftmargin}{30pt}\setlength{\parsep}{0pt}\usecounter{spacerule}}
\item Problems 7 and 8 in the previous week's assignment.
\item OpenStax Ch.~6, problems 65, 70, 86, 103, 105, 139, 146, 148, 163, 168.
\item Strang \S 15.1, problems 15, 32. 
\item Strang \S 15.2, problems 22, 28.  
\item Strang \S 15.3, problems 29, 33.
\end{list}
\end{assignment}
\begin{assignment}({\em{\color{duecolor} For the week ending on February 26, 2020}})
\begin{list}{\textbf{\arabic{spacerule})}}{\setlength{\itemsep}{4pt}\setlength{\leftmargin}{30pt}\setlength{\parsep}{0pt}\usecounter{spacerule}}
\item OpenStax Ch.~6, problems 328, 335, 338, 341, 364, 367, 375.
\item Strang \S 15.6, problems 10, 12, 16, 17, 28, 30, 46. 

\end{list}
\end{assignment}
\begin{assignment}({\em{\color{duecolor} For the week ending on March 4, 2020}})
\begin{list}{\textbf{\arabic{spacerule})}}{\setlength{\itemsep}{4pt}\setlength{\leftmargin}{30pt}\setlength{\parsep}{0pt}\usecounter{spacerule}}
\item OpenStax Ch.~6, problems 385, 388, 394, 406, 408, 424.
\item Strang \S 15.5, problems 12, 15, 18, 25.
\item Lax, Terrell, problems 8.26, 8.28. 
\end{assignment}
\begin{assignment}({\em{\color{duecolor} For the week ending on March 11, 2020}})
\begin{list}{\textbf{\arabic{spacerule})}}{\setlength{\itemsep}{4pt}\setlength{\leftmargin}{30pt}\setlength{\parsep}{0pt}\usecounter{spacerule}}
\item Lax, Terrell, problems 9.4, 9.5 b),c), 9.6 b), d), 9.10, 9.11. 
\end{list}
\end{assigment}
\end{document}

%%%%%%%%%%%%%%%%%%%%%%%%%%%%%%%%%%%%%%%%%%%%%%%%%%%%%
\begin{assignment}({\em{\color{duecolor} Due October 14, 2019}})
\begin{list}{\textbf{\arabic{spacerule})}}{\setlength{\itemsep}{2pt}\setlength{\leftmargin}{30pt}\setlength{\parsep}{0pt}\usecounter{spacerule}}
\item R\&R, problem 4, $(a),\dots(d)$, page 39.
\item R\&R, problem 8, page 47.
\item R\&R, problem 3, page 59.
\item R\&R, problem 22 (b), page 60.
\end{list}
\end{assignment}

\begin{assignment}({\em{\color{duecolor} Due October 21, 2019}})
\begin{list}{\textbf{\arabic{spacerule})}}{\setlength{\itemsep}{2pt}\setlength{\leftmargin}{30pt}\setlength{\parsep}{0pt}\usecounter{spacerule}}
\item\begin{enumerate}[(a)]
\item Show that for all $n\geq5$, $2^n>n^2$.
\vskip2pt
\item Prove that for all $n\geq1$, 
\begin{equation*}
1-\dfrac12+\dfrac13-\dfrac14+\cdots+\dfrac1{2n-1}-\dfrac1{2n}=
\dfrac1{n+1}+\dfrac1{n+2}+\cdots+\dfrac1{2n}\,.
\end{equation*}
\end{enumerate}
\item R\&R, problem 13, page 70.
\item R\&R, problem 10, page 76.
\end{list}

\end{assignment}

\begin{assignment}
({\em{\color{duecolor} Due October 28, 2019}})
\newline
\indent\textbf{0)} Problem 3, homework assignment III.
\begin{list}{\textbf{\arabic{spacerule})}}{\setlength{\itemsep}{2pt}\setlength{\leftmargin}{30pt}\setlength{\parsep}{0pt}\usecounter{spacerule}}
\item R\&R, problem 8, page 76 (You need not consider a leap year).
\item R\&R, problem 14, page 76.
\end{list}
\end{assignment}
\begin{assignment}({\em{\color{duecolor} Due November 13, 2019}})
\begin{list}{\textbf{\arabic{spacerule})}}{\setlength{\itemsep}{2pt}\setlength{\leftmargin}{30pt}\setlength{\parsep}{0pt}\usecounter{spacerule}}
\item R\&R, problem 22, page 87.
\item R\&R, problem 1, $(b)$, $(e)$, page 94.
\item R\&R, problem 8, page 95.
\end{list}
\end{assignment}

\begin{assignment}({\em{\color{duecolor} Due November 25, 2019}})
\begin{list}{\textbf{\arabic{spacerule})}}{\setlength{\itemsep}{2pt}\setlength{\leftmargin}{30pt}\setlength{\parsep}{0pt}\usecounter{spacerule}}
\item R\&R, problem 11, page 101
\item R\&R, problem 15, page 101 
\item R\&R, problem 6, page 109
\end{list}
\end{assignment}

\begin{assignment}({\em{\color{duecolor} Due December 4, 2019}})
\begin{list}{\textbf{\arabic{spacerule})}}{\setlength{\itemsep}{2pt}\setlength{\leftmargin}{30pt}\setlength{\parsep}{0pt}\usecounter{spacerule}}
\item R\&R, problem 6, $(b)$, $(c)$, $(d)$, page 131.
\item R\&R, problem 6, page 140.
\item R\&R, problem 7, $(c)$, page 146.
\end{list}
\end{assignment}
\end{document}
\begin{problem}
Let $X$ be a set and let $\beta$ be a collection of subsets of $X$ such that 
\begin{list}{\textbf{\roman{spacerule})}}{\setlength{\itemsep}{2pt}\setlength{\leftmargin}{30pt}\setlength{\parsep}{0pt}\usecounter{spacerule}}
\item For every $p\in X$ there is some $B\in\beta$ such that $p\in B$.
\item Let $p\in B_1\cap B_2$, where $B_1, B_2\in\beta$. Then there is some $B_3\in\beta$ such that $p\in B_3\subset B_1\cap B_2$.
\end{list}
Show that the family of subsets of $X$ consisting of all unions of the members of $\beta$ together with the empty set $\emptyset$ form a topology on $X$. Explain why the above statement is stronger (and hence more useful) that Theorem 2.5 in the course text.
\end{problem}
\begin{problem} Show that the following collections $\tau$ of subsets of $X$ form a topology in the given space.
\begin{list}{\textbf{\alph{spacerule})}}{\setlength{\itemsep}{2pt}\setlength{\leftmargin}{30pt}\setlength{\parsep}{0pt}\usecounter{spacerule}}
\item Let $X=\R{}$ with $\tau$ consisting of all subsets $B$ of $\R{}$ such that $\R{}\setminus B$ contains finitely many elements or is all of $\R{}$.
\begin{comment}
\item Let $X = [-1,1]$ with $\tau$ consisting of the sets $[-1,b)$, $(a,1]$, $(a,b)$, where $0<b<1$, $-1<a<0$, together with $\emptyset$ and $[-1,1]$.
\end{comment}
\item Let $X=\{a,b,c\}$ and $\tau = \{\emptyset, \{c\}, \{a,c\}, \{b,c\}, \{a,b,c\}\}$.
\end{list}
\end{problem}
\begin{problem}
Let $X=\R{}$ and suppose that $\R{}$ is equipped with the topology defined by the basis $\beta=\{[a,b)\,\vert\,a<b\}$. Call this space $\R{}_\tau$. Which of the following functions $f\colon \R{}_\tau\to\R{}_\tau$ are continuous?
\begin{list}{\textbf{\alph{spacerule})}}{\setlength{\itemsep}{2pt}\setlength{\leftmargin}{30pt}\setlength{\parsep}{0pt}\usecounter{spacerule}}
\item $f(x) = 2x-5$
\item $f(x) = -x$
\item $f(x) = x^2$
\end{list}
\end{problem}
\end{document}
\begin{problem}
Let $\openset\subset\R2$ be an open set. A \emph{proper coordinate patch} is a one-to-one immersion $\mathbf{x}\colon\openset\to\R3$ such that the inverse function $\mathbf{x}^{-1}\colon\mathbf{x}(\openset)\to\openset$ is continuous in the subspace topology of $\mathbf{x}(\openset)$. A \emph{surface} is a subset $\smfld\subset\R3$ equipped with the subspace topology such that for each point $\pnt\in\smfld$, there is a proper coordinate patch whose image contains a neighborhood of $\pnt$ in $\smfld{}$. 
\begin{list}{\textbf{\alph{spacerule})}}{\setlength{\itemsep}{2pt}\setlength{\leftmargin}{30pt}\setlength{\parsep}{0pt}\usecounter{spacerule}}
\item The torus $T^2$ is the image of the mapping \[\mathbf{x}(u,v) = ((R+r\cos u)\cos v,(R+r\cos u)\sin v,r\sin u),\qquad r<R.\] Show that the torus is a surface.
\item Prove that a surface is also a differentiable manifold.  
\end{list}
\end{problem}
\newcounter{projective}
\begin{problem}\setcounter{projective}{\theprob}
The complex projective space $\CP1$ is defined as the set of all $1$ dimensional complex subspaces of $\cx2=\cx{}\times\cx{}$, or alternatively, the orbit space of the action of $\cxp{}=\cx{}\smallsetminus\{0\}$ on $\cxp2=\cx2\smallsetminus\{(0,0)\}$ by multiplication. Write $\projP\colon\cxp2\to\CP1$ for the projection. Endow $\CP1$ with the usual quotient topology so that a set $U\subset\CP1$ is open if and only if the union of all lines constituting $U$ (with the origin removed) is open in $\cxp2$, cf.~Lee, pp.~604-606 and \href{http://en.wikipedia.org/wiki/Quotient_space_%28topology%29}{\textbf{\color{hyperlink}Quotient space}} in Wikipedia.
\newcounter{atlas} 
\newcounter{charts}
\setcounter{spacerule}{1}
\begin{list}{\textbf{\alph{spacerule})}}{\setlength{\itemsep}{2pt}\setlength{\leftmargin}{30pt}\setlength{\parsep}{0pt}\usecounter{spacerule}}
\item Explain why the quotient topology on $\CP1$ is second countable and Hausdorff, and show that the projection map $\projP$ is open.
\item\setcounter{charts}{\thespacerule} Define an atlas of coordinate charts on $\CP1$ in analogy with the standard coordinate charts for $\RP1$ constructed in class, and show that these are homeomorphisms. Also identify the image of each coordinate map. 
\item \setcounter{atlas}{\thespacerule} Find the transition functions ({\it a.k.a.~}the change of coordinate maps) for the atlas for $\CP1$ constructed in part \alph{charts}. 
\end{list}
\end{problem}

\begin{problem}
Let $\CP1$ be the complex projective space as in problem \Roman{projective}.
\setcounter{spacerule}{1}
\begin{list}{\textbf{\alph{spacerule})}}{\setlength{\itemsep}{2pt}\setlength{\leftmargin}{30pt}\setlength{\parsep}{0pt}\usecounter{spacerule}}
\item Show that the projection $\projP\colon\cxp2\to\CP1$ is a smooth mapping.
\item Compute the rank of the projection $\projP$ at every point in $\cxp2$.
\item Use part \alph{atlas} of problem \Roman{projective} to show that $\CP1$ is diffeomorphic to the unit sphere $S^2$.
\end{list}
\end{problem}
\end{document}
%\begin{problem}
%A regular space curve $\curve$ is called a (generalized) helix if there is a constant unit vector $\unitv$ so that the angle between $\T(s)$ and $\unitv$ does not depend on $s$. Show that a space curve with $\curvature(s)>0$ is a helix if and only if $\torsion(s)/\curvature(s)$ is constant.  
%\end{problem}
\newcounter{bases}
\newcounter{grassmann}
\begin{problem}\setcounter{grassmann}{\theprob}
Let $\Gr n 2$ denote the Grassmannian space of $2$\,--\,dimensional (vector) subspaces in $\R n$. 
\setcounter{spacerule}{1}
\begin{list}{\textbf{\alph{spacerule})}}{\setlength{\itemsep}{2pt}\setlength{\leftmargin}{30pt}\setlength{\parsep}{0pt}\usecounter{spacerule}}
\item\setcounter{bases}{\thespacerule} Let $\plane\in\Gr n 2$. Then any two bases $\{\vb_1,\vb_2\}$, $\{\wb_1,\wb_2\}$ for $\plane$ are related by 
\begin{equation}\label{eq:equivalence}
\wb_i = \sum_{j=1,2} a_i^j\vb_j,\qquad i=1,2,
\end{equation}
where $(a_i^j)\in GL(2)$ is an invertible $2\times 2$ matrix. Conclude that $\Gr n 2$
can be identified with the equivalence classes of linearly independent pairs of vectors $\{\vb_1,\vb_2\}$, where $\{\vb_1,\vb_2\}\sim\{\wb_1,\wb_2\}$ provided that equation \eqref{eq:equivalence} holds for some $(a_i^j)\in GL(2)$.
\item Let $GL(2)$ act on the space $\rktwos$ of $2\times n$ matrices of rank $2$ by left multiplication. Use part \alph{bases} to show that the orbit space $\rktwos/GL(2)$ of the action can be identified with $\Gr n 2$.
\item Equip $\rktwos$ with the subspace topology as an open subset of $\R{2n}$, and equip $\Gr n 2$ with the usual quotient topology. Conclude that the projection $\projGr\colon\rktwos\to\Gr n 2$ is an open mapping.  
\item Let $\rktwo i j\subset\rktwos$ denote the set of rank $2$ matrices whose minor consisting of the $i$th and $j$th columns is invertible. Show that the action of $GL(2)$ on $\rktwos$ preserves the sets $\rktwo i j$ and that every $GL(2)$ orbit in $\rktwo i j$ contains a unique matrix whose $i,j$-minor is the identity matrix.  
\item\setcounter{atlas}{\thespacerule} Conclude that every plane $\plane\in\rktwo i j$ can be identified with a unique $(n-2)\times2$ matrix. (For example, with $n=4$ and $i=1$, $j=3$, the $2\times2$ matrix $\left(\begin{matrix} a&b\\c&d\end{matrix}\right)$ would correspond to the plane admitting the basis $\vb_1=(1,a,0,b)$, $\vb_2=(0,c,1,d)$.)
\item Show that each $\Grs i j=\projGr(\rktwo i j)\subset\Gr n 2$ is open, and use part \alph{atlas} to define coordinate maps $\varphi_{ij}$ in each $\Grs i j$. In particular, show that each $\varphi_{ij}\colon\Grs i j\to \R{2n-4}$ is a homeomorphism.
\item Finally compute a representative sample of transition functions to conclude that $\Gr n 2$ forms a differentiable manifold. 
\end{list}
\end{problem}

\begin{problem} Let $\CP n$ and $\Gr n 2$ denote the complex projective space and Grassmann manifold constructed in problems $\theprojective$ and $\thegrassmann$.
\setcounter{spacerule}{1}
\begin{list}{\textbf{\alph{spacerule})}}{\setlength{\itemsep}{2pt}\setlength{\leftmargin}{30pt}\setlength{\parsep}{0pt}\usecounter{spacerule}}
\item Let $S^3$ be the unit sphere in $\cx2$ identified with $\R4$, and let the Hopf map $\pi_{H}:S^3\to S^2$ be the restriction of the projection $\projP$ to $S^3$. Describe the inverse image $\pi_H{}^{-1}(p)$ of a point $p\in S^2$.  

\item Show that the projection $\projGr\colon\rktwos\to\Gr n 2$ is differentiable.
\item Show that each $\Gr n 2$ is compact.
\item Let $\vsp\subset\R n$, where $n\geq4$, be a fixed $4$-dimensional vector space, and let $\smfld\subset\Gr n 2$ consist of all planes $\plane$ contained in $\vsp$. Is $\smfld$ an embedded submanifold of $\Gr n 2$? 

\end{list}
\end{problem}
\end{document}
\begin{problem}
Find all unit speed curves $\curve(s)$ with $\curvature(s)=1/(1+s^2)$ and $\torsion(s)=0$.
\end{problem}

\begin{problem}
Let $\curve(t)$, $\curveb(t)$ be two regular curves defined on some interval $I$. Then $\curveb$ is called an {\em involute} of $\curve$ provided that for all $t\in I$, $\curveb(t)$ lies on the tangent line to $\curve$ at $\curve(t)$, and the respective tangent vectors to $\curve$ and $\curveb$ at $\curve(t)$ and $\curveb(t)$ are perpendicular.
\setcounter{spacerule}{1}
\begin{list}{\textbf{\alph{spacerule})}}{\setlength{\itemsep}{2pt}\setlength{\parsep}{0pt}\setlength{\leftmargin}{32pt}\usecounter{spacerule}}
\item Suppose that $\curve(s)$ is unit speed. Then prove that an involute $\curveb$ of $\curve$ is given by $\curveb(s) = \curve(s)+(c-s)\T(s)$, where $c$ is a constant and $T(s)=\curve'(s)$. Interpret the expression for $\curveb$ geometrically.
\item Prove that the involute of a slope line is a plane curve.
\end{list}
\end{problem}
\end{document}
